% =========================================================================
% 六、问题二的结果分析
% =========================================================================

\section{问题二的结果分析}

结果分析按``总体—分日—风险—对照''四个层次展开：先给出评价期总体结果（\S\ref{sec:q2-total}），再分析指定日期的调度细节（\S\ref{sec:q2-days}），继而讨论全年运行形态与逐月费用（\S\ref{sec:q2-year}）以及紧急购电与风险结构（\S\ref{sec:q2-risk}），最后通过对照方案与消融实验定位各项机制的价值（\S\ref{sec:q2-ablation}）。

\subsection{评价期总费用与构成}
\label{sec:q2-total}

\begin{table}[H]
  \centering
  \caption{主方案评价期结果（2025.2.1--12.31，334 天，起评 10800 kWh）}
  \label{tab:q2-total}
  \begin{tabular}{lc}
    \toprule
    项目 & 数值 \\
    \midrule
    \textbf{总费用} & \textbf{1395.6995 万元} \\
    计划费用 & 1318.2416 万元（94.5\%） \\
    紧急费用 & 77.4578 万元（5.5\%） \\
    计划购电量 & 2165.26 万 kWh \\
    紧急购电量 & 12.44 万 kWh \\
    发生紧急购电的天数 & 163 / 334（48.8\%） \\
    弃置供能余量 & 252.63 万 kWh \\
    \bottomrule
  \end{tabular}
\end{table}

紧急费用只占总费用的 5.5\%，紧急电量仅占计划购电量的 0.57\%。这说明星期偏差校正已修正预测的中心趋势，分位裕量（$\beta=0.75$）足以吸收绝大部分残余波动，5 倍电价惩罚仅在少数时段被触发。

\subsection{指定日期的调度结果}
\label{sec:q2-days}

\begin{figure}[H]
  \centering
  \includegraphics[width=0.86\textwidth]{p2_fig3.pdf}
  \caption{指定日期的调度细节：(a)~分时电价与日前计划购电量；(b)~净负载实测、预测与紧急购电；(c)~储能日内充放电动作；(d)~储能电量轨迹}
  \label{fig:q2-days}
\end{figure}

下列数值按字面时钟区间报告；\texttt{result2.xlsx} 按官方模板原始行序逐行填写、未改动模板标签（模板行标签较字面区间整体后移一个区间），两口径相差一个 10 分钟区间，核对时需注意。

\begin{table}[H]
  \centering
  \caption{指定日期的计划购电量（对应题目表 1 的指定时间段，kWh）}
  \label{tab:q2-t1}
  \small
  \setlength{\tabcolsep}{5pt}
  \begin{tabular}{|c|c|c|c|c|}
    \hline
    时间段 & 2025-03-20 & 2025-06-21 & 2025-09-23 & 2025-12-21 \\ \hline
    10:00--10:10 & 0.0000 & 0.0000 & 0.0000 & 0.0000 \\ \hline
    12:00--12:10 & 545.0194 & 0.0000 & 406.6138 & 1042.1389 \\ \hline
    14:00--14:10 & 0.0000 & 0.0000 & 0.0000 & 0.0000 \\ \hline
    16:00--16:10 & 512.0462 & 224.7427 & 499.3924 & 852.9497 \\ \hline
    18:00--18:10 & 701.8570 & 0.0000 & 772.4848 & 735.2130 \\ \hline
    20:00--20:10 & 0.0000 & 0.0000 & 0.0000 & 0.0000 \\ \hline
  \end{tabular}
\end{table}

\begin{table}[H]
  \centering
  \caption{指定日期的全天购电量（kWh）与费用汇总（对应题目表 1 的汇总行）}
  \label{tab:q2-daysum}
  \small
  \setlength{\tabcolsep}{5pt}
  \begin{tabular}{|c|r|r|r|r|}
    \hline
    日期 & 计划购电量 & 紧急购电量 & 合计购电量 & 费用 (元) \\ \hline
    2025-03-20 & 67872.0855 & 293.0319 & 68165.1174 & 43077.5075 \\ \hline
    2025-06-21 & 36887.8821 & 0.0000 & 36887.8821 & 20795.3006 \\ \hline
    2025-09-23 & 65531.4200 & 1317.0219 & 66848.4419 & 49169.8748 \\ \hline
    2025-12-21 & 98705.7213 & 0.0000 & 98705.7213 & 63439.8523 \\ \hline
  \end{tabular}
\end{table}

表中``全天计划购电量''为 $b_t$ 全天之和，\textbf{不含}紧急购电量；``合计购电量''为两者之和；全天费用＝计划费＋紧急费（式~\eqref{eq:q2-cost}）。

\begin{table}[H]
  \centering
  \caption{指定日期储能设备充放电量（kWh）与日初日末储电量（对应题目表 2）}
  \label{tab:q2-t2}
  \small
  \setlength{\tabcolsep}{4pt}
  \begin{tabular}{|c|c|c|c|c|c|}
    \hline
    时间段 & 充电量 & 放电量 & 时间段 & 充电量 & 放电量 \\ \hline
    \multicolumn{6}{|c|}{2025-03-20} \\ \hline
    0:00--4:00   & 133.3  & 189.8   & 4:00--8:00   & 300.1   & 5166.4 \\ \hline
    8:00--12:00  & 6059.0 & 2777.4  & 12:00--16:00 & 3905.3  & 543.2 \\ \hline
    16:00--20:00 & 234.1  & 5937.8  & 20:00--24:00 & 10717.5 & 2692.8 \\ \hline
    \multicolumn{2}{|c|}{0:00 储电量} & 10800.0000
      & \multicolumn{2}{c|}{24:00 储电量} & 10783.9968 \\ \hline
    \multicolumn{6}{|c|}{2025-06-21} \\ \hline
    0:00--4:00   & 169.8   & 3413.0  & 4:00--8:00   & 384.1   & 4836.6 \\ \hline
    8:00--12:00  & 9630.8  & 0.0000  & 12:00--16:00 & 0.0000  & 0.0000 \\ \hline
    16:00--20:00 & 115.0   & 5727.6  & 20:00--24:00 & 10025.0 & 2485.9 \\ \hline
    \multicolumn{2}{|c|}{0:00 储电量} & 10800.0000
      & \multicolumn{2}{c|}{24:00 储电量} & 10800.0000 \\ \hline
    \multicolumn{6}{|c|}{2025-09-23} \\ \hline
    0:00--4:00   & 330.5  & 288.5   & 4:00--8:00   & 411.2   & 6550.3 \\ \hline
    8:00--12:00  & 4827.1 & 1896.6  & 12:00--16:00 & 5621.7  & 596.0 \\ \hline
    16:00--20:00 & 35.0   & 6159.9  & 20:00--24:00 & 10298.8 & 2061.4 \\ \hline
    \multicolumn{2}{|c|}{0:00 储电量} & 10585.6180
      & \multicolumn{2}{c|}{24:00 储电量} & 10454.3513 \\ \hline
    \multicolumn{6}{|c|}{2025-12-21} \\ \hline
    0:00--4:00   & 94.5   & 172.0   & 4:00--8:00   & 989.4   & 1528.6 \\ \hline
    8:00--12:00  & 6188.6 & 6906.2  & 12:00--16:00 & 6116.1  & 2237.9 \\ \hline
    16:00--20:00 & 0.0000 & 5712.0  & 20:00--24:00 & 10561.2 & 2842.5 \\ \hline
    \multicolumn{2}{|c|}{0:00 储电量} & 10800.0000
      & \multicolumn{2}{c|}{24:00 储电量} & 10800.0000 \\ \hline
  \end{tabular}
\end{table}

四日都呈现``白天光伏充沛时充电、晚峰放电、深夜低谷再充满''的典型日循环形态（图~\ref{fig:q2-days}c、d）；日末与日初电量并不相等（如 09-23 日初 10585.62 $\to$ 日末 10454.35 kWh），体现了实际储电量的跨日连续衔接。

\begin{table}[H]
  \centering
  \caption{指定日期紧急购电结果（对应题目表 3，合并连续区间）}
  \label{tab:q2-t3}
  \small
  \setlength{\tabcolsep}{6pt}
  \begin{tabular}{|c|c|r|}
    \hline
    日期 & 紧急购电时间段 & 紧急购电量 (kWh) \\ \hline
    2025-03-20 & 20:30--20:40 & 293.03 \\ \hline
    2025-06-21 & （无紧急购电） & 0.00 \\ \hline
    2025-09-23 & 20:20--20:40 & 1126.01 \\ \hline
               & 20:50--21:10 & 129.61 \\ \hline
               & 21:20--21:50 & 61.40 \\ \hline
    2025-12-21 & （无紧急购电） & 0.00 \\ \hline
  \end{tabular}
\end{table}

紧急购电集中在晚间光伏退出后的时段（09-23 共 1317.02 kWh，见图~\ref{fig:q2-days}b），与净负载的日内形态一致；06-21 为光伏旺季，计划自足；12-21 为冬季大负载日，计划量充裕。

\subsection{全年运行形态与逐月费用}
\label{sec:q2-year}

\begin{figure}[H]
  \centering
  \includegraphics[width=0.9\textwidth]{p2_fig4.pdf}
  \caption{逐日购电费用构成与储能电量全年轨迹}
  \label{fig:q2-year}
\end{figure}

\begin{figure}[H]
  \centering
  \includegraphics[width=0.88\textwidth]{p2_fig5.pdf}
  \caption{月度费用构成与累计费用}
  \label{fig:q2-month}
\end{figure}

\begin{table}[H]
  \centering
  \caption{主方案月度费用（万元）与风险概况}
  \label{tab:q2-month}
  \begin{tabular}{crcrrr}
    \toprule
    月份 & 总费用 & 计划 & 紧急 & 大紧急天数$^{*}$ & 触顶天数 \\
    \midrule
    2  & 125.96 & 124.23 & 1.74 & 0 & 28 \\
    3  & 117.46 & 107.58 & 9.88 & 3 & 31 \\
    \textbf{4}  & \textbf{98.84}（最低） & 91.75 & 7.09 & 3 & 30 \\
    5  & 99.13 & 92.45 & 6.68 & 6 & 31 \\
    6  & 144.34 & 126.52 & 17.83 & 5 & 28 \\
    7  & 150.79 & 143.39 & 7.40 & 4 & 30 \\
    8  & 130.64 & 125.27 & 5.37 & 4 & 31 \\
    9  & 114.24 & 108.07 & 6.16 & 4 & 28 \\
    10 & 117.45 & 109.96 & 7.49 & 5 & 31 \\
    11 & 126.43 & 124.86 & 1.57 & 0 & 30 \\
    \textbf{12} & \textbf{170.42}（最高） & 164.16 & 6.26 & 2 & 29 \\
    \bottomrule
  \end{tabular}
\end{table}

$^{*}$\ 大紧急天数＝当日紧急购电量超过 1000 kWh 的天数。

12 月费用最高（170.42 万元，冬季大负载、光伏弱，计划购电量最多）；4 月最低（98.84 万元，春季温和、光伏进入旺季）；6 月紧急费用最高（17.83 万元，5 个大紧急日），对应初夏光伏出力波动叠加晚峰缺口，是全年风险最集中的月份（图~\ref{fig:q2-month}）。

\subsection{紧急购电与风险结构}
\label{sec:q2-risk}

\begin{figure}[H]
  \centering
  \includegraphics[width=0.9\textwidth]{p2_fig6.pdf}
  \caption{紧急购电特征：(a)~紧急购电量时空分布；(b)~日紧急购电费用分布}
  \label{fig:q2-emerg}
\end{figure}

\begin{table}[H]
  \centering
  \caption{主方案风险指标（评价期 334 天）}
  \label{tab:q2-risk}
  \begin{tabular}{lc}
    \toprule
    指标 & 数值 \\
    \midrule
    紧急费用分位数 P50 / P90 / P95 / P99 & 0.00 / 0.71 / 1.33 / 3.20 万元 \\
    单日最坏紧急费用 & 7.58 万元 \\
    大紧急天数（$>1000$ kWh） & 36 / 334（10.8\%） \\
    触底天数（日内有 $E_t\le E_{\min}$） & 113 / 334（33.8\%） \\
    触顶天数（日内有 $E_t=E_{\max}$） & 327 / 334（97.9\%） \\
    日均起始／日内最低／日内最高储电量 & 10673 / \textbf{1444} / \textbf{10793} kWh \\
    年末储电量（12.31 日末） & 10667 kWh \\
    \bottomrule
  \end{tabular}
\end{table}

\textbf{储备结构的准确表述}：过半数日期（P50＝0）不发生任何紧急购电，尾部风险可控（P99＝3.20 万元，图~\ref{fig:q2-emerg}）。触顶天数 327/334 仅表示这些天日内\textbf{至少一次}达到上限，不等于全天维持满电：日内最低储电量平均 1444 kWh，接近下限 1200 kWh，说明储能每天仍经历一次幅度较大的充放电循环（日均最高 10793 $\to$ 日均最低 1444 kWh）。因此主方案的运行形态是``低谷充满—晚峰放空—再充满''的日循环（图~\ref{fig:q2-year}），而非``长期保持满电''。

\subsection{方案对照与条件消融}
\label{sec:q2-ablation}

本文在建模过程中先后构造了一组递进方案，用于对照不同的预测器与执行规则。它们与主方案共用评价口径（2/1 起评 10800 kWh），但分属不同的模型族，故\textbf{仅作水平对照，不将差额归因于某一单项机制}（图~\ref{fig:q2-ladder}a）。

\begin{table}[H]
  \centering
  \caption{对照方案与主方案的水平比较（评价口径统一）}
  \label{tab:q2-ladder}
  \small
  \begin{tabular}{llcc}
    \toprule
    方案 & 预测方式／执行方式 & 总费用 (万元) & 相对 A \\
    \midrule
    A & 固定代表日曲线／贪心 & 2825.28 & — \\
    B & 代表日＋残差中位数校准／贪心 & 2035.24 & $-28.0\%$ \\
    B48 & B＋48 小时计划窗口／贪心 & 2024.98 & $-28.3\%$ \\
    C & B＋分位裕量（$\beta=0.7$）／贪心 & 1775.91 & $-37.1\%$ \\
    D & C／每 4 小时重优化 & 1739.64 & $-38.4\%$ \\
    \textbf{E（主方案）} & \textbf{近期水平＋星期偏差＋分位裕量}／\textbf{尽限充放电} & \textbf{1395.70} & \textbf{$-50.6\%$} \\
    完美信息参照 & 当天实测（不可实施）／尽限充放电 & 1229.10 & $-56.5\%$ \\
    \bottomrule
  \end{tabular}
\end{table}

\begin{figure}[H]
  \centering
  \includegraphics[width=0.9\textwidth]{p2_fig7.pdf}
  \caption{(a)~各方案的评价期总费用阶梯；(b)~条件消融实验}
  \label{fig:q2-ladder}
\end{figure}

\textbf{同链路的单项消融。}为单独识别星期偏差校正的作用，另做``关闭星期偏差''的同链路对照（图~\ref{fig:q2-ladder}b）：保持计划层、执行层、起评电量、$\beta=0.70$、$W=14$ 全部不变，仅令 $\gamma=0$，得 1719.71 万元。与同一链路下 $N_{\mathrm{DOW}}=4$ 的主方案相比，星期偏差校正与逐月联合选参合计带来 \textbf{324.01 万元（$-18.84\%$）}的降幅。该对照中各项设置一致，差额可归因于预测器与选参方式；但两项各自的贡献尚未单独分离。

\textbf{与完美信息参照的差距。}主方案 1395.70 万元，完美信息参照 1229.10 万元，相差 166.60 万元。这一差距不应整体解读为``可改进空间''：① 参照方案同样受日循环约束，若取消该约束其费用只会更低，故 166.60 万元不是严格意义上的信息缺口成本；② 差距中包含天气等不可预测的随机成分，任何因果策略都必须承担这部分代价；③ 5 倍电价的损失不对称使``保险溢价''不可避免——即使预测无偏，最优计划水平也应高于条件中位数。
